
\documentclass{article}
\usepackage{colortbl}
\usepackage{makecell}
\usepackage{multirow}
\usepackage{supertabular}

\begin{document}

\newcounter{utterance}

\centering \large Interaction Transcript for game `dond', experiment `semi', episode 0 with nemotron{-}3{-}ultra.
\vspace{24pt}

{ \footnotesize  \setcounter{utterance}{1}
\setlength{\tabcolsep}{0pt}
\begin{supertabular}{c@{$\;$}|p{.15\linewidth}@{}p{.15\linewidth}p{.15\linewidth}p{.15\linewidth}p{.15\linewidth}p{.15\linewidth}}
    \# & $\;$A & \multicolumn{4}{c}{Game Master} & $\;\:$B\\
    \hline

    \theutterance \stepcounter{utterance}  
    & & \multicolumn{4}{p{0.6\linewidth}}{
        \cellcolor[rgb]{0.9,0.9,0.9}{
            \makecell[{{p{\linewidth}}}]{
                \texttt{\tiny{[P1$\langle$GM]}}
                \texttt{Egy kooperatív tárgyalási játékban játszol, amelyben meg kell egyeznöd egy másik játékossal arról, hogyan osztozz meg egy tárgyak halmazán közöttetek.} \\
\\ 
\texttt{Szabályok:} \\
\texttt{(a) Te és a másik játékos tárgyak halmazát kapod. Mindannyian kapsz egy titkos értékfüggvényt is, amely azt mutatja, hogy mennyire értékelsz minden tárgy típust.} \\
\texttt{(b) Üzeneteket váltasz a másik játékossal, hogy megállapodj abban, hogy ki mit kap. Mindegyik játékos legfeljebb 5 üzenetet küldhet, vagy bármikor korábban befejezheti a játékot egy titkos javaslat meghozatalával.} \\
\texttt{(c) Egy titkos javaslatot kell benyújtanod, amely megjelöli a kívánt tárgyakat, négyzet zárójelben formázva az alábbiak szerint: "{[}Javaslat: <szám> <tárgy neve>, <szám> <tárgy neve>, <...>{]}"} \\
\texttt{(d) Ha javaslataid kiegészítők, azaz elegendő tárgy van mind a két javaslat teljesítéséhez, minden játékos pontszámot kap az általa kapott tárgyak értékeinek összege alapján. Különben mindketten nulla pontot kapsz.} \\
\texttt{(e) A céld az, hogy maximalizáld a pontszámod és a másik játékos pontszámának összegét.} \\
\\ 
\texttt{Kezdjük el.} \\
\\ 
\texttt{Az elérhető tárgyak halmaza:} \\
\\ 
\texttt{1 kanál, 1 szék, 2 virágok, 2 lámpák} \\
\\ 
\texttt{A titkos értékfüggvényed:} \\
\\ 
\texttt{kanál: 2, szék: 2, virág: 0, lámpa: 3} \\
\\ 
\texttt{FONTOS: Üzeneteid, ha nem titkos javaslat, közvetlenül átadódnak a másik játékosnak, ezért ne válaszolj a szabályokra vagy olyan szövegre, amely üzenetedről szól. Titkos javaslat meghozatalához használd a megadott formátumot. Ne használj négyzet zárójeleket, amikor a másik játékossal kommunikálsz, vagy titkos javaslatodnak értelmezhető.} \\
            }
        }
    }
    & \\ \\

    \theutterance \stepcounter{utterance}  
    & \multicolumn{4}{p{0.6\linewidth}}{
        \cellcolor[rgb]{0.9,0.9,0.9}{
            \makecell[{{p{\linewidth}}}]{
                \texttt{\tiny{[P1$\rangle$GM]}}
                \texttt{Szia! Láthatod, hogy 2 lámpa, 1 szék és 1 kanál van rendelkezésre. Nekem a lámpák értékük 3, a szék és a kanál pedig 2. A virágok számomra 0 értékűek.} \\
\\ 
\texttt{Milyen az értékfüggvényed? Ha neked a virágok, a szék vagy a kanál a magasabb értékű, akkor könnyen megoldhatjuk az elosztást.} \\
            }
        }
    }
    & & \\ \\

    \theutterance \stepcounter{utterance}  
    & & & \multicolumn{2}{p{0.3\linewidth}}{
        \cellcolor[rgb]{0.9,0.9,0.9}{
            \makecell[{{p{\linewidth}}}]{
                \texttt{\tiny{[GM$|$GM]}}
                \texttt{continue} \\
            }
        }
    }
    & & \\ \\

    \theutterance \stepcounter{utterance}  
    & & & \multicolumn{2}{p{0.3\linewidth}}{
        \cellcolor[rgb]{0.9,0.9,0.9}{
            \makecell[{{p{\linewidth}}}]{
                \texttt{\tiny{[GM$|$GM]}}
                \texttt{Szia! Láthatod, hogy 2 lámpa, 1 szék és 1 kanál van rendelkezésre. Nekem a lámpák értékük 3, a szék és a kanál pedig 2. A virágok számomra 0 értékűek.} \\
\\ 
\texttt{Milyen az értékfüggvényed? Ha neked a virágok, a szék vagy a kanál a magasabb értékű, akkor könnyen megoldhatjuk az elosztást.} \\
            }
        }
    }
    & & \\ \\

    \theutterance \stepcounter{utterance}  
    & & \multicolumn{4}{p{0.6\linewidth}}{
        \cellcolor[rgb]{0.9,0.9,0.9}{
            \makecell[{{p{\linewidth}}}]{
                \texttt{\tiny{[GM$\rangle$P2]}}
                \texttt{Egy kooperatív tárgyalási játékban játszol, amelyben meg kell egyeznöd egy másik játékossal arról, hogyan osztozz meg egy tárgyak halmazán közöttetek.} \\
\\ 
\texttt{Szabályok:} \\
\texttt{(a) Te és a másik játékos tárgyak halmazát kapod. Mindannyian kapsz egy titkos értékfüggvényt is, amely azt mutatja, hogy mennyire értékelsz minden tárgy típust.} \\
\texttt{(b) Üzeneteket váltasz a másik játékossal, hogy megállapodj abban, hogy ki mit kap. Mindegyik játékos legfeljebb 5 üzenetet küldhet, vagy bármikor korábban befejezheti a játékot egy titkos javaslat meghozatalával.} \\
\texttt{(c) Egy titkos javaslatot kell benyújtanod, amely megjelöli a kívánt tárgyakat, négyzet zárójelben formázva az alábbiak szerint: "{[}Javaslat: <szám> <tárgy neve>, <szám> <tárgy neve>, <...>{]}"} \\
\texttt{(d) Ha javaslataid kiegészítők, azaz elegendő tárgy van mind a két javaslat teljesítéséhez, minden játékos pontszámot kap az általa kapott tárgyak értékeinek összege alapján. Különben mindketten nulla pontot kapsz.} \\
\texttt{(e) A céld az, hogy maximalizáld a pontszámod és a másik játékos pontszámának összegét.} \\
\\ 
\texttt{Kezdjük el.} \\
\\ 
\texttt{Az elérhető tárgyak halmaza:} \\
\\ 
\texttt{1 kanál, 1 szék, 2 virágok, 2 lámpák} \\
\\ 
\texttt{A titkos értékfüggvényed:} \\
\\ 
\texttt{kanál: 3, szék: 1, virág: 3, lámpa: 0} \\
\\ 
\texttt{FONTOS: Üzeneteid, ha nem titkos javaslat, közvetlenül átadódnak a másik játékosnak, ezért ne válaszolj a szabályokra vagy olyan szövegre, amely üzenetedről szól. Titkos javaslat meghozatalához használd a megadott formátumot. Ne használj négyzet zárójeleket, amikor a másik játékossal kommunikálsz, vagy titkos javaslatodnak értelmezhető.} \\
\\ 
\texttt{Szia! Láthatod, hogy 2 lámpa, 1 szék és 1 kanál van rendelkezésre. Nekem a lámpák értékük 3, a szék és a kanál pedig 2. A virágok számomra 0 értékűek.} \\
\\ 
\texttt{Milyen az értékfüggvényed? Ha neked a virágok, a szék vagy a kanál a magasabb értékű, akkor könnyen megoldhatjuk az elosztást.} \\
            }
        }
    }
    & \\ \\

    \theutterance \stepcounter{utterance}  
    & & & \multicolumn{4}{p{0.6\linewidth}}{
        \cellcolor[rgb]{0.9,0.9,0.9}{
            \makecell[{{p{\linewidth}}}]{
                \texttt{\tiny{[GM$\langle$P2]}}
                \texttt{Szia! Nekem a virágok és a kanál a magas értékűek (mindkettő 3 pont), a szék 1 pont, a lámpák pedig 0 pont.} \\
\\ 
\texttt{Úgy tűnik, a legjobban úgy járunk, ha a virágokat és a kanált én kapom, a lámpákat és a széket pedig te. Így mindkettőnk 6 pontot kap.} \\
\\ 
\texttt{Milyen ez a javaslat számodra?} \\
            }
        }
    }
     \\ \\

    \theutterance \stepcounter{utterance}  
    & & & \multicolumn{2}{p{0.3\linewidth}}{
        \cellcolor[rgb]{0.9,0.9,0.9}{
            \makecell[{{p{\linewidth}}}]{
                \texttt{\tiny{[GM$|$GM]}}
                \texttt{continue} \\
            }
        }
    }
    & & \\ \\

    \theutterance \stepcounter{utterance}  
    & & & \multicolumn{2}{p{0.3\linewidth}}{
        \cellcolor[rgb]{0.9,0.9,0.9}{
            \makecell[{{p{\linewidth}}}]{
                \texttt{\tiny{[GM$|$GM]}}
                \texttt{Szia! Nekem a virágok és a kanál a magas értékűek (mindkettő 3 pont), a szék 1 pont, a lámpák pedig 0 pont.} \\
\\ 
\texttt{Úgy tűnik, a legjobban úgy járunk, ha a virágokat és a kanált én kapom, a lámpákat és a széket pedig te. Így mindkettőnk 6 pontot kap.} \\
\\ 
\texttt{Milyen ez a javaslat számodra?} \\
            }
        }
    }
    & & \\ \\

    \theutterance \stepcounter{utterance}  
    & & \multicolumn{4}{p{0.6\linewidth}}{
        \cellcolor[rgb]{0.9,0.9,0.9}{
            \makecell[{{p{\linewidth}}}]{
                \texttt{\tiny{[P1$\langle$GM]}}
                \texttt{Szia! Nekem a virágok és a kanál a magas értékűek (mindkettő 3 pont), a szék 1 pont, a lámpák pedig 0 pont.} \\
\\ 
\texttt{Úgy tűnik, a legjobban úgy járunk, ha a virágokat és a kanált én kapom, a lámpákat és a széket pedig te. Így mindkettőnk 6 pontot kap.} \\
\\ 
\texttt{Milyen ez a javaslat számodra?} \\
            }
        }
    }
    & \\ \\

    \theutterance \stepcounter{utterance}  
    & \multicolumn{4}{p{0.6\linewidth}}{
        \cellcolor[rgb]{0.9,0.9,0.9}{
            \makecell[{{p{\linewidth}}}]{
                \texttt{\tiny{[P1$\rangle$GM]}}
                \texttt{Ez a javaslat tökéletesen elfogadható számomra! Így kapom a két lámpát (2×3=6 pont) és a széket (2 pont), összesen 8 pontot, te meg a két virágot (2×3=6 pont) és a kanált (2 pont), összesen 8 pontot. Ez maximális az összegben is.} \\
\\ 
\texttt{{[}Javaslat: 2 lámpa, 1 szék{]}} \\
            }
        }
    }
    & & \\ \\

    \theutterance \stepcounter{utterance}  
    & & & \multicolumn{2}{p{0.3\linewidth}}{
        \cellcolor[rgb]{0.9,0.9,0.9}{
            \makecell[{{p{\linewidth}}}]{
                \texttt{\tiny{[GM$|$GM]}}
                \texttt{proposal} \\
            }
        }
    }
    & & \\ \\

    \theutterance \stepcounter{utterance}  
    & & & \multicolumn{2}{p{0.3\linewidth}}{
        \cellcolor[rgb]{0.9,0.9,0.9}{
            \makecell[{{p{\linewidth}}}]{
                \texttt{\tiny{[GM$|$GM]}}
                [0, 1, 0, 2]
            }
        }
    }
    & & \\ \\

    \theutterance \stepcounter{utterance}  
    & & \multicolumn{4}{p{0.6\linewidth}}{
        \cellcolor[rgb]{0.9,0.9,0.9}{
            \makecell[{{p{\linewidth}}}]{
                \texttt{\tiny{[GM$\rangle$P2]}}
                \texttt{A másik játékos benyújtott egy titkos javaslatot. Most rajtad a sor, hogy te is benyújtsd a saját javaslatodat.} \\
            }
        }
    }
    & \\ \\

    \theutterance \stepcounter{utterance}  
    & & & \multicolumn{4}{p{0.6\linewidth}}{
        \cellcolor[rgb]{0.9,0.9,0.9}{
            \makecell[{{p{\linewidth}}}]{
                \texttt{\tiny{[GM$\langle$P2]}}
                \texttt{{[}Javaslat: 1 kanál, 2 virágok{]}} \\
            }
        }
    }
     \\ \\

    \theutterance \stepcounter{utterance}  
    & & & \multicolumn{2}{p{0.3\linewidth}}{
        \cellcolor[rgb]{0.9,0.9,0.9}{
            \makecell[{{p{\linewidth}}}]{
                \texttt{\tiny{[GM$|$GM]}}
                \texttt{proposal} \\
            }
        }
    }
    & & \\ \\

    \theutterance \stepcounter{utterance}  
    & & & \multicolumn{2}{p{0.3\linewidth}}{
        \cellcolor[rgb]{0.9,0.9,0.9}{
            \makecell[{{p{\linewidth}}}]{
                \texttt{\tiny{[GM$|$GM]}}
                [1, 0, 2, 0]
            }
        }
    }
    & & \\ \\

    \theutterance \stepcounter{utterance}  
    & & & \multicolumn{2}{p{0.3\linewidth}}{
        \cellcolor[rgb]{0.9,0.9,0.9}{
            \makecell[{{p{\linewidth}}}]{
                \texttt{\tiny{[GM$|$GM]}}
                [[0, 1, 0, 2], [1, 0, 2, 0]]
            }
        }
    }
    & & \\ \\

\end{supertabular}
}

\end{document}
